% from main.tex
%%%%%%%%%%  scan idx 621  |  folio 608  %%%%%%%%%%
\textbf{Corollaire 3.4.} Soient \(k\) un corps, \(X\) un \(k\)-schéma propre. On suppose que \(\Pic_{X_{\red}/k}\) possède un faisceau inversible universel. (Par exemple, ce sera le cas, d'après la théorie de descente, si toute composante connexe \(X'_i\) de \(X_{\red}\) possède un point rationnel, par exemple si \(k\) est algébriquement clos, ou plus généralement si le p.g.c.d. des degrés des \(0\)-cycles de tout \(X'_i\) est \(1\)). Alors au-dessus de tout ouvert affine \(V\) de \(\Pic_{X/k}\), il existe un faisceau inversible universel sur \(X\times_kV\).

Finalement, on démontre I sous la forme plus précise suivante~:

\textbf{Proposition 3.5.} Soient \(S\) un schéma noethérien intègre, \(f:X\hookrightarrow Y\) une immersion nilpotente de \(S\)-schémas propres. Alors il existe un ouvert non vide \(V\) de \(S\) tel que le morphisme
\[
  f^*:\Pic_{Y/S}|V\longrightarrow\Pic_{X/S}|V
\]
soit représentable par un morphisme affine de présentation finie~; en fait, \(N=\Ker f^*\) est un schéma en groupes affines à fibres unipotentes (au sens de SGA 3 XVII).

En effet \(X\) est défini par un idéal nilpotent \(I\). Alors \(I\) possède une filtration par des idéaux \(I_0=I\supset I_1\supset\cdots\supset I_n=0\) tels que \(NI_i\subset I_{i-1}\), où \(N\) est le nilradical de \(\mathcal O_Y\) (par exemple \(I_k=N^kI\)). Soit \(X_i\) le sous-schéma de \(Y\) défini par \(I_i\). Par récurrence il suffit d'étudier les immersions \(X_i\hookrightarrow X_{i-1}\). On peut donc supposer que \(N.I=0\)~; c'est la condition i) de 3.1 et on finit en se ramenant au cas où les autres conditions (ii) à (v) sont vérifiées également. Grâce au théorème de la platitude générique (EGA IV 6.9.1) appliqué à \(\mathcal O_{X'}\), et aux \(R^if_*(I)\) et à \(E\), et compte tenu des relations de Künneth EGA III 6.9.9, on peut supposer, quitte à restreindre \(S\), que les conditions ii), iii) et iv) sont vérifiées.

